\subsection{\texorpdfstring{应用：\(R^{n}\) 中欧氏球的测度}{应用：R to the n 中欧氏球的测度}}

设 \(\mu\) 是 \(\mathbf{R}\) 上的 Lebesgue 测度，\(\mu_n\) 是 \(\mathbf{R}^n\) 上的 Lebesgue 测度，即 \(n\) 个都等于 \(\mu\) 的因子之积。我们来计算欧氏单位球的测度 \(V_n = \mu_n(\mathbf{B}_n)\) 。由命题 7 的推论 2，

\[
\mathrm{V} _ {n} = \int_ {- 1} ^ {+ 1} \mu_ {n - 1} \big (\mathbf {B} _ {n} (z _ {n}) \big) d z _ {n}.\tag{16}
\]

现在，截口 \(\mathbf{B}_n(z_n)\) 是 \(\mathbf{R}^{n-1}\) 中由关系式 \(\sum_{i=1}^{n-1} z_i^2 \leqslant 1 - z_n^2\) 定义的子集，换句话说，它是球 \(\mathbf{B}_{n-1}\) 在比为 \(\sqrt{1 - z_n^2}\) 的位似下的像。但由命题 11 及公式

\[
\alpha \int_ {- \infty} ^ {+ \infty} f (\alpha x) d x = \int_ {- \infty} ^ {+ \infty} f (z) d z
\]

（对 \(f \in \mathcal{K}(\mathbf{R})\) ）立即得出，\(\mu_{n-1}\) 在位似 \(\mathbf{x} \mapsto \alpha \mathbf{x}\) 下的像是测度 \(\alpha^{1-n} \mu_{n-1}\) 。因此

\[
\mu_ {n - 1} \left(\mathbf {B} _ {n} \left(z _ {n}\right)\right) = \left(\sqrt {1 - z _ {n} ^ {2}}\right) ^ {n - 1} \mathrm{V} _ {n - 1}.
\]

把它代入 (16)，并作变量替换 \(z_{n} = \sin \varphi\) （其中 \(-\frac{\pi}{2} \leqslant \varphi \leqslant \frac{\pi}{2}\) ），即得

\[
\mathrm{V} _ {n} = \mathrm{V} _ {n - 1} \int_ {- \frac {\pi}{2}} ^ {+ \frac {\pi}{2}} \cos^ {n} \varphi d \varphi = 2 \mathrm{V} _ {n - 1} \int_ {0} ^ {\frac {\pi}{2}} \cos^ {n} \varphi d \varphi .\tag{17}
\]

但（FRV, Ch. VII, §1, No.~3, 公式 (20)）

\[
\int_ {0} ^ {\frac {\pi}{2}} \cos^ {m} \varphi d \varphi = \frac {1}{2} \frac {\Gamma \left(\frac {1}{2}\right) \Gamma \left(\frac {m + 1}{2}\right)}{\Gamma \left(\frac {m + 2}{2}\right)}
\]

而把上式代入关系式 (17) 并注意到 \(\Gamma(\frac{1}{2})\) 的表达式（FRV, VII, §1, No.~3, 公式 (21)），最终得到

\[
\mathrm{V} _ {n} = \frac {\pi^ {n / 2}}{\Gamma \left(\frac {n}{2} + 1\right)}.\tag{18}
\]

